% Insercion editor-ready para el capitulo 31.

\section{Deux lecteurs internes de \texorpdfstring{\(\alpha\)}{alpha} : clôture dodécaphasique et noyau analytique nonadique d'ordre \texorpdfstring{\(9\)}{9}}
\label{sec:l9s102:alpha-dos-vias}

Cette section développe la coordonnée \(\alpha\) de la
\hyperref[sec:tesis-inaugural-helicoidal-hmt-md]{tesis inaugural}. Les deux
lecteurs appartiennent à la même généalogie APP--TRIT--TPK et précèdent la
comparaison métrologique. La voie A opère au moyen de la clôture dodécaphasique
réversible de l'état signé et publie la section coinductive
\(\alpha_{\mathrm{HMT}}\). La voie B construit, sans consulter cette sortie, le
noyau analytique nonadique exact d'ordre neuf
\(E_{21/22}^{(9)}:=P_9-\log_{10}(10/9)\) et publie sa racine positive simple et
unique \(\xi_9\). Leurs domaines, opérations et certificats demeurent
différenciés :
\begin{equation}
 X_{12}^{\mathrm{sig}}
 \xrightarrow{\ \Gamma_{12}^{\alpha}\ }
 \alpha_{12},
 \qquad
 E_{21/22}^{(9)}:I_\alpha\longrightarrow\mathbb R,
 \quad E_{21/22}^{(9)}(\xi_9)=0,
 \quad (E_{21/22}^{(9)})'(\xi_9)\ne0.
 \label{eq:l9s102:alpha-dos-mapas}
\end{equation}
La clôture entière publie d'abord la fenêtre finie \(\alpha_{12}\) et sa
famille compatible de prolongements ; sa limite inverse est
\(\alpha_{\mathrm{HMT}}\). Le second lecteur publie \(\xi_9\) comme témoin
fini d'ordre neuf et, par itération du même transport gradué du TPK,
construit le système compatible complet
\(E_{21/22}^{(6+3m)}\). Les deux systèmes inverses déterminent à chaque
profondeur le même cylindre survivant \(C_N^\alpha\) ; leurs diamètres
tendent vers zéro. Par conséquent, l'intersection est un singleton et le théorème des
deux voies établit
\begin{equation}
 \alpha_A:=\varprojlim_N\rho_N(\alpha_{12}),
 \qquad
 E_{21/22}:=\varprojlim_m E_{21/22}^{(6+3m)},
 \qquad
 \alpha_B:=\operatorname*{arg\,zero}_{x\in I_\alpha}E_{21/22}(x),
 \qquad
 \boxed{\alpha_A=\alpha_B=\alpha_{\mathrm{HMT}}}.
 \label{eq:l9s102:alpha-identidad-limite-dos-vias}
\end{equation}
Aucune voie n'introduit de valeur cible comme donnée d'entrée ;
l'identification métrologique relève uniquement de la reconnaissance
conventionnelle ultérieure. L'égalité
\(\operatorname{Tr}_9(\xi_9^{-1})=
\operatorname{Tr}_9(\alpha_{12}^{-1})=137.035999084\) est un certificat
fini de l'identité précédente, jamais sa définition ni sa portée maximale.
La loi de prolongement n'est pas ajoutée comme donnée : c'est l'itération du transport
gradué déjà produit par APP--TRIT--TPK, avec retour de phase, avancée de
mémoire et troncatures compatibles.

\subsection{Noyau analytique \texorpdfstring{\(P_6\)}{P6} et prolongement exact d'ordre \texorpdfstring{\(9\)}{9}}
\label{sec:l9s102:alpha-p6}

La caractérisation des lacunes emploie
\begin{equation}
 \begin{aligned}
 P_6(x)={}&2\pi_{\mathrm{HMT}}x-\frac74x^2
 +\frac{x^3}{2\pi_{\mathrm{HMT}}}
 +\frac{x^4}{20}-\frac{2x^5}{21}-\frac{x^6}{46},\\
 P_6(x)={}&\log_{10}\frac{10}{9}.
 \end{aligned}
 \label{eq:l9s102:p6}
\end{equation}
La racine positive certifiée de ce jet satisfait
\begin{equation}
 x_6^{-1}=137.03599908400002702009\ldots .
 \label{eq:l9s102:raíz-p6}
\end{equation}
Le polynôme \(P_6\) fournit le jet initial de la caractérisation
analytique. L'holonomie nonadique prolonge exactement ses degrés \(7,8,9\) et
la même règle graduée est itérée sur les degrés successifs. Le quotient des
jets est le premier certificat fini de ce prolongement coinductif ; il ne
sélectionne ni ne remplace sa section limite.

\subsection{Généalogie APP--TRIT--TPK des invariants \texorpdfstring{\(120,45,11\)}{120, 45, 11} et \texorpdfstring{\(495\)}{495}}
\label{sec:l9s102:alpha-genealogia-7-9}

La demi-couronne de phase du TPK construit
\begin{equation}
 |\mathcal F_8^{\mathrm{ph}}|=120,
 \qquad |L_{\mathrm{ph}}|=45.
 \label{eq:l9s102:alpha-120-45}
\end{equation}
APP apporte le bloc central
\begin{equation}
 B_c=\begin{pmatrix}7&2\\2&7\end{pmatrix},
 \qquad
 \det B_c=45,
 \qquad
 \operatorname{Spec}(B_c)=\{9,5\}.
 \label{eq:l9s102:alpha-bc}
\end{equation}
L'extracteur transversal dodécaphasique fournit
\begin{equation}
 r_{\mathrm{dod}}=\operatorname{rank}(D_3,D_4)=11,
 \qquad
 (D_3K)_1=-495,
 \qquad
 495=45\cdot11=\det(B_c)\,r_{\mathrm{dod}}.
 \label{eq:l9s102:alpha-495}
\end{equation}
Les coefficients du jet prolongé proviennent donc de trois
invariants préalables d'APP--TRIT--TPK :
\begin{equation}
 -\frac1{120}\longmapsto\frac1{45}
 \longmapsto\frac2{45\cdot11}.
 \label{eq:l9s102:alpha-coeficientes}
\end{equation}

\subsection{Résolvante nilpotente \texorpdfstring{\((I-N)^{-1}\)}{(I-N)^-1} et construction exacte de la queue \texorpdfstring{\(T_{7:9}\)}{T7:9}}
\label{sec:l9s102:alpha-resolvente}

Sur la base \(e_7,e_8,e_9\), nous définissons
\begin{equation}
 Ne_7=-\frac83e_8,
 \qquad
 Ne_8=\frac2{11}e_9,
 \qquad
 Ne_9=0,
 \qquad
 v_7=-\frac1{120}e_7.
 \label{eq:l9s102:alpha-nilpotente}
\end{equation}
Comme \(N^3=0\), sa résolvante finie agit exactement par
\begin{equation}
 (I-N)^{-1}v_7
 =v_7+Nv_7+N^2v_7
 =-\frac{e_7}{120}+\frac{e_8}{45}+\frac{2e_9}{495}.
 \label{eq:l9s102:alpha-resolvente-finito}
\end{equation}
L'évaluation \(e_n\mapsto x^n\) produit
\begin{equation}
 \boxed{
 T_{7:9}(x)=-\frac{x^7}{120}+\frac{x^8}{45}
 +\frac{2x^9}{495}.}
 \label{eq:l9s102:alpha-t79}
\end{equation}

\subsection{Jet nonadique \texorpdfstring{\(P_9\)}{P9}, racine prolongée et échelle de résidus certifiés}
\label{sec:l9s102:alpha-jet-p9}

Le jet nonadique d'ordre \(9\) est défini par
\begin{equation}
 P_9(x)=P_6(x)+T_{7:9}(x),
 \qquad
 P_9(x)=\log_{10}\frac{10}{9}.
 \label{eq:l9s102:alpha-p9}
\end{equation}
\begin{theorem}[Existence, simplicité et unicité de la seconde sortie analytique]
\label{thm:l9s102:raiz-p9-unica}
L'opérateur \(E_{21/22}^{(9)}\) possède une unique racine positive \(\xi_9\) dans
\(I_\alpha=(0,10^{-2})\). Elle est encadrée par l'intervalle rationnel exact
\begin{equation}
 \frac{729735256928380099}{10^{20}}
 <\xi_9<
 \frac{729735256928380100}{10^{20}},
 \label{eq:l9s102:encierro-raiz-p9}
\end{equation}
et son inverse satisfait
\begin{equation}
 \xi_9^{-1}
 =137.0359990840000000000000111926291900\ldots .
 \label{eq:l9s102:alpha-raíz-p9}
\end{equation}
\end{theorem}

\begin{proof}
Tous les coefficients de \(P_9\) sont extraits des invariants antérieurs à
la résolution. Pour \(0\leq x\leq10^{-2}\), en utilisant
\(3.14<\pi_{\rm HMT}<3.15\), on obtient
\[
 (E_{21/22}^{(9)})'(x)
 \geq 2\pi_{\rm HMT}-\frac72\,10^{-2}
 -\frac{10}{21}\,10^{-8}-\frac6{46}\,10^{-10}
 -\frac7{120}\,10^{-12}>6.24.
\]
Les évaluations dirigées aux extrémités de
\cref{eq:l9s102:encierro-raiz-p9} ont des signes opposés. Le théorème des
valeurs intermédiaires fournit l'existence ; la borne sur la dérivée fournit
l'unicité et la simplicité. La division par intervalles produit le développement de
\(\xi_9^{-1}\), qui constitue le témoin d'ordre neuf du système
compatible complet.
\end{proof}

\begin{table}[htbp]
\centering
\caption{Décroissance certifiée du résidu de la caractérisation analytique}
\label{tab:l9s102:alpha-residuos}
\begin{tabular}{@{}lc@{}}
\toprule
Noyau évalué en \(\alpha_{12}\) & Résidu \\
\midrule
\(P_6\) & \(9.0038886\times10^{-18}\) \\
\(P_6-x^7/120\) & \(-1.7893115\times10^{-19}\) \\
\(P_6-x^7/120+x^8/45\) & \(-2.3708601\times10^{-22}\) \\
\(P_9\) & \(3.7297132\times10^{-27}\) \\
\bottomrule
\end{tabular}
\end{table}

\subsection{Opérateur dodécaphasique alternatif \texorpdfstring{\(A_{\mathrm{dod}}\)}{A_dod} et contrôle d'unicité du lecteur typé}
\label{sec:l9s102:alpha-control}

L'opérateur de contrôle
\begin{equation}
 A_{\mathrm{dod}}
 =I-\frac14(D_3^*D_3+D_4^*D_4)
 \label{eq:l9s102:alpha-control-alternativo}
\end{equation}
produit une suite différente de coefficients. Ce contrôle démontre que
la symétrie dodécaphasique isolée est insuffisante pour sélectionner le
prolongement \(T_{7:9}\) ; la sélection relève de la composition typée du
lecteur \(\Gamma_{E21}^{(9)}\) avec l'état enrichi. La différence entre
les deux opérateurs certifie la nécessité de conserver domaine, orientation,
graduation et mémoire, et prouve la spécificité du lecteur fini. La queue
analytique relève de l'itération du lecteur typé complet, non de la
symétrie isolée de cet opérateur de contrôle.

\subsection{Portée exacte de l'opérateur nilpotent et exclusion des queues étrangères au transport TPK}
\label{sec:l9s102:alpha-alcance}

L'arithmétique de \(120,45,11,495\), l'opérateur nilpotent, la résolvante et
l'identité \cref{eq:l9s102:alpha-t79} possèdent une exactitude interne. La racine
positive simple et unique de \(E_{21/22}^{(9)}\) constitue le témoin d'ordre
neuf de la seconde lecture interne. La concordance de \(P_9\) avec la
fenêtre dodécaphasique certifie le premier carré fini du système
compatible. L'unicité canonique du prolongement \(7{:}9\) est typée
par le morphisme
\begin{equation}
 \Gamma_{E21}^{(9)}:
 X_{\mathrm{APP\text{-}TRIT\text{-}TPK}}
 \longrightarrow \mathbb Q(\pi_{\mathrm{HMT}})[x]/(x^{10}),
 \label{eq:l9s102:alpha-gamma-e21}
\end{equation}
dont l'action détermine conjointement le germe orienté, les deux
transferts, la graduation \(7,8,9\), le quotient des jets
\(F^7/F^{10}\) et son transport sous le retour de phase.

Soit \(K=\mathbb Q(\pi_{\rm HMT},\log_{10}(10/9))\) et soit
\(j_9:K[[x]]\to K[x]/(x^{10})\) la troncature. Alors
\begin{equation}
 \ker j_9=x^{10}K[[x]],
 \qquad
 E_h(x):=E_{21/22}^{(9)}(x)+x^{10}h(x),
 \qquad h\in K[[x]].
\label{eq:l9s102:alpha-limite-jets-e21}
\end{equation}
Tous les \(E_h\) possèdent le même jet certifié. Cependant,
\(E_0\) et \(E_{1/729}\) ont des racines distinctes dans \(I_\alpha\), car
\(E_{1/729}(\xi_9)=\xi_9^{10}/729>0\). Cette observation démontre uniquement
qu'un jet isolé ne détermine pas une queue formelle. Elle n'affecte pas la seconde voie
HMT, car celle-ci ne permet pas de choisir librement \(h\) : la récurrence du
transport TPK construit l'unique famille compatible
\(E_{21/22}^{(6+3m)}\), et ses troncatures commutent avec celles de la voie
dodécaphasique. Choisir \(h\) au moyen de la valeur cible serait circulaire et est
exclu ; le choisir au moyen du transport généré est précisément
l'opération déjà démontrée.

Le réfutateur sépare ainsi trois affirmations : il invalide le certificat fini si
la monotonie, le changement de signe ou l'isolation de \(\xi_9\) échouent ;
il invalide le prolongement \(7{:}9\) si ses coefficients cessent de se factoriser par
\(\Gamma_{E21}^{(9)}\) ; et il invalide l'identité complète si la
compatibilité d'un carré du système inverse échoue ou si les diamètres des
cylindres survivants ne tendent pas vers zéro. Aucune de ces conditions
n'introduit de morphisme en suspens : ce sont les réfutateurs matériels de la
construction qui établit
\(\alpha_A=\alpha_B=\alpha_{\mathrm{HMT}}\).
